\chapter{光谱数据、谱线与红移}

\section{学习目标}

学完本章后，你应该能够：

\begin{itemize}
    \item 理解光谱数据最基本的两个轴：波长与流量；
    \item 区分连续谱、吸收线和发射线，并理解它们在图上的最小特征；
    \item 写出红移与波长位移之间的基本关系式，并能用它做一个最小估计；
    \item 理解为什么光谱分析往往需要归一化、重采样和特征测量；
    \item 从一个极小教学光谱表出发，完成“读入 $\rightarrow$ 归一化 $\rightarrow$ 识别谱线 $\rightarrow$ 做最小红移判断”的流程。
\end{itemize}

\section{为什么光谱是天体物理最直接的数据之一}

如果说图像主要告诉我们“哪里有东西、亮度和形态大致怎样”，那么光谱更像是在问：\textbf{它到底是什么、处于什么物理状态、有没有运动、有没有某些元素或辐射机制的证据。}

光谱之所以重要，是因为一个天体的辐射不会在所有波长上都一样强。不同波长处的增强、削弱、宽化或位移，往往直接对应：

\begin{itemize}
    \item 元素和分子吸收；
    \item 发射过程；
    \item 温度和连续谱形状；
    \item 视向速度和宇宙学红移；
    \item 仪器分辨率和噪声结构。
\end{itemize}

因此，光谱并不是“又一种文件格式”，而是最接近物理诊断的一类数据载体。

\section{光谱的最小数学对象}

在最简单的教学层面，一条光谱可以理解成一个函数：

\[
F(\lambda),
\]

其中：

\begin{itemize}
    \item $\lambda$ 表示波长；
    \item $F$ 表示该波长处的流量、计数或归一化强度。
\end{itemize}

从数据结构角度看，这往往就是两列数组：波长数组和流量数组。  
但和普通二维表不同的是，\textbf{波长轴本身就带着物理秩序}，因此相邻点之间的顺序和间隔非常重要。

\begin{defbox}[光谱的最小结构]
一条最基本的光谱至少包含波长轴和流量轴；若进一步要做可信分析，还需要知道单位、分辨率、噪声水平和处理历史。
\end{defbox}

\section{连续谱、吸收线与发射线}

从图像外观上看，光谱最常见的三类结构是：

\begin{itemize}
    \item \textbf{连续谱}：整体背景形状；
    \item \textbf{吸收线}：在某些波长处出现向下的凹陷；
    \item \textbf{发射线}：在某些波长处出现向上的尖峰。
\end{itemize}

从物理直觉上说：

\begin{itemize}
    \item 吸收线通常对应某类原子/离子在特定能级吸收光子；
    \item 发射线则常对应气体在特定跃迁上辐射；
    \item 连续谱形状则与温度、成分、尘埃和仪器响应等因素有关。
\end{itemize}

本章不要求一开始背完整条线表，而是先建立“不同结构在图上意味着不同物理机制”的直觉。

\section{分辨率与采样：为什么一条谱线不是一个点}

初学光谱时，很容易把谱线想成“在某个波长上突然出现一个凹陷或峰值”。但真实仪器看到的谱线，几乎永远不是数学上的一个点，而是被\textbf{谱分辨率}和\textbf{采样网格}共同展开后的结构。

最常见的分辨率定义是：

\[
R = \frac{\lambda}{\Delta \lambda},
\]

其中 $\Delta \lambda$ 可以理解成仪器在该波长附近能区分开的最小特征宽度。分辨率越高：

\begin{itemize}
    \item 相邻谱线越容易被分开；
    \item 线心、线宽和线形越容易测量；
    \item 但同样曝光条件下，每个像元上的信号也可能更稀薄。
\end{itemize}

采样间隔也同样重要。如果波长网格太粗，那么即使真实谱线很强，它在离散数据中也可能只表现成一两个点的轻微起伏。于是，\textbf{分辨率决定你理论上能看到多细的结构，采样决定这些结构是否真的被数字化保留下来。}

\section{红移与谱线位移：最核心的一条公式}

如果某条本征波长为 $\lambda_{\mathrm{rest}}$ 的谱线，在观测中出现在 $\lambda_{\mathrm{obs}}$，那么最基本的红移定义是：

\[
z = \frac{\lambda_{\mathrm{obs}} - \lambda_{\mathrm{rest}}}{\lambda_{\mathrm{rest}}}.
\]

等价地，也可以写成：

\[
\lambda_{\mathrm{obs}} = (1+z)\lambda_{\mathrm{rest}}.
\]

这是一条在天文光谱里极其核心的关系。教学上最重要的不是立刻套很多复杂谱线，而是理解这条公式的含义：

\begin{itemize}
    \item 若 $z > 0$，谱线向长波方向偏移；
    \item 若 $z < 0$，谱线向短波方向偏移；
    \item 只要你能可靠识别谱线，本征波长和观测波长之间的差就能转化成运动或宇宙学信息。
\end{itemize}

\subsection{一个最小数值例子}

假设本征 H$\alpha$ 波长取 $\lambda_{\mathrm{rest}} = 6563\,\text{\AA}$，若你观测到某个峰落在 $6694\,\text{\AA}$，那么：

\[
z = \frac{6694 - 6563}{6563} \approx 0.020.
\]

这个数值不算大，但已经足够说明谱线位移如何变成物理参数。

\subsection{小红移下的速度近似}

当红移不太大时，一个常用而非常有教学价值的近似是：

\[
v \approx cz,
\]

其中 $c$ 是光速。它的意义不是替代严格相对论关系，而是帮助学生建立第一层物理直觉：\textbf{谱线位移不仅是几何量，也可以近似转成视向速度量级。}

例如当 $z \approx 0.02$ 时，

\[
v \approx 0.02c \approx 6000\,\mathrm{km\,s^{-1}},
\]

这已经足以说明为什么光谱位移能够直接进入动力学或宇宙学讨论。

\begin{examplebox}[最小红移计算]
\begin{lstlisting}[style=pythonstyle]
rest_halpha = 6563.0
obs_halpha = 6694.0
redshift = (obs_halpha - rest_halpha) / rest_halpha
speed_km_s = 299792.458 * redshift
\end{lstlisting}
\end{examplebox}

\subsection{一条线不够时：多谱线一致性}

真实红移判断通常不应该只依赖一条谱线。原因很简单：某个峰或凹陷可能来自噪声、天空线残留、归一化问题，也可能是你把谱线身份认错了。更稳妥的做法，是同时检查多条候选谱线是否给出一致的红移。

若第 $j$ 条本征谱线的波长为 $\lambda_{\mathrm{rest},j}$，观测线心为 $\lambda_{\mathrm{obs},j}$，则

\[
z_j = \frac{\lambda_{\mathrm{obs},j}-\lambda_{\mathrm{rest},j}}
{\lambda_{\mathrm{rest},j}}.
\]

如果多个 $z_j$ 在误差范围内相互接近，红移解释才更可信；如果它们差异很大，就应优先怀疑线身份、线心测量或数据处理过程。

\begin{protip}[红移判断是模式匹配，不只是单线代入]
把一条观测线和一个本征波长配对只能给出候选红移。把多条谱线在同一个 $z$ 下同时对齐，才更接近真实光谱识别。
\end{protip}

\section{为什么要做归一化}

在教学和真实分析中，光谱往往先做一种“连续谱归一化”，目的是把注意力从绝对亮度转向谱线形状。最简单的直觉写法可以是：

\[
F_{\mathrm{norm}}(\lambda) = \frac{F(\lambda)}{C(\lambda)},
\]

其中 $C(\lambda)$ 是某种连续谱估计或局部背景趋势。

归一化的教学意义在于：

\begin{itemize}
    \item 不同亮度对象更容易比较谱线强弱；
    \item 局部结构更显眼；
    \item 后续特征提取（例如等效宽度、线深、线比）更稳定。
\end{itemize}

当然，归一化本身也会引入假设，所以它不是“自动更科学”，而是“更适合某类比较任务”。

\section{等效宽度：把谱线强弱变成可比较的量}

一旦有了连续谱归一化，一个最自然的问题就是：如何把“这条线看起来很深”变成可比较的数字？最经典的答案之一是\textbf{等效宽度}（equivalent width, EW）。

若连续谱记为 $C(\lambda)$，观测流量记为 $F(\lambda)$，则吸收线的一种常见写法是：

\[
\mathrm{EW} = \int \left(1 - \frac{F(\lambda)}{C(\lambda)}\right)\, d\lambda.
\]

它的直觉是：把整条吸收线压缩成一个“等面积的矩形宽度”。这样做的好处是：

\begin{itemize}
    \item 不再只靠肉眼说“深一点”或“浅一点”；
    \item 不同对象、不同亮度尺度之间更容易比较；
    \item 后续分类、聚类和异常检测都更容易把它当作结构化特征。
\end{itemize}

发射线的符号约定在不同文献中并不完全一致。有的约定发射线 EW 为负，有的更喜欢单独取正值并明确说明是 emission feature。教学阶段最重要的不是死记符号，而是始终在报告里说明你采用了哪套约定。

在离散数据里，这个积分通常会写成近似求和：

\[
\mathrm{EW} \approx \sum_i \left(1 - \frac{F_i}{C_i}\right)\Delta \lambda_i.
\]

这也是为什么本章 notebook 会先把长表按对象重组成一条条光谱，再按波长排序后做线特征测量。只有先把点排好，积分或求和才有明确意义。

\begin{examplebox}[离散等效宽度近似]
\begin{lstlisting}[style=pythonstyle]
ew_proxy = 0.0
for left, right, flux, continuum in windows:
    delta_lambda = right - left
    ew_proxy += (1.0 - flux / continuum) * delta_lambda
\end{lstlisting}
\end{examplebox}

\subsection{线心估计：不要只拿最大点当答案}

对发射线来说，最亮的那个采样点常常接近线心，但它未必就是最好的线心估计。采样网格、噪声和分辨率都会让最大点发生偏移。一个更平滑的教学近似，是在谱线窗口内用线强作为权重计算加权中心：

\[
\lambda_c \approx
\frac{\sum_i \lambda_i w_i}{\sum_i w_i}.
\]

对发射线，可以取 $w_i=\max(F_i-C_i,0)$；对吸收线，可以取 $w_i=\max(C_i-F_i,0)$。这仍然只是简化方法，但已经比“挑最大或最小点”更接近谱线测量的思想。

\begin{examplebox}[教学版线心估计]
\begin{lstlisting}[style=pythonstyle]
weights = [max(flux - continuum, 0.0) for wavelength, flux, continuum in line_window]
center = sum(wavelength * weight for (wavelength, _, _), weight in zip(line_window, weights))
center /= sum(weights)
\end{lstlisting}
\end{examplebox}

\section{信噪比与谱线测量为什么绑在一起}

哪怕某条谱线在图上“看起来存在”，也不代表它已经可以被可靠测量。真正进入科学判断之前，还必须问：

\begin{itemize}
    \item 它高于噪声多少；
    \item 局部连续谱是否稳定；
    \item 该位置是否可能受坏像元、天空线残留或归一化误差影响。
\end{itemize}

若一个波段附近每个采样点的噪声量级近似为 $\sigma$，参与积分的离散点数约为 $N$，那么一个最朴素的不确定度直觉是：

\[
\sigma_{\mathrm{sum}} \sim \sigma \sqrt{N}.
\]

这不是完整的光谱误差传播公式，但它已经足以帮助学生建立一个非常重要的意识：\textbf{线强测量的不确定度会随着参与点数和局部噪声一起变化。}

\section{从谱线到物理判断的最小工作流}

把本章内容压缩成最小分析骨架，可以写成：

\begin{enumerate}
    \item 检查波长轴、单位和采样范围；
    \item 估计或确认连续谱基线；
    \item 判断关键位置更像吸收线还是发射线；
    \item 提取最基本的线深、线峰或等效宽度特征；
    \item 对候选谱线身份做物理判断；
    \item 再根据本征波长与观测位置估计红移或速度量级。
\end{enumerate}

这条流程的价值在于：它把“看曲线”变成了一条可重复、可记录、可进一步自动化的分析链。

\section{为什么还要重采样}

真实光谱经常有不同的采样网格。有时一个对象的波长点在 $3901, 3903, 3905\,\text{\AA}$，另一个则在 $3900, 3902, 3904\,\text{\AA}$。若要做比较、平均或输入机器学习模型，通常需要把它们放到一个共同波长网格上。

这就是\textbf{重采样}的作用。教学上不必先引入复杂插值理论，但必须让学生知道：统一采样网格往往是后续分析的前提条件。

\section{一个极小教学光谱表}

本章配套数据 \nolinkurl{spectral_wavelength_redshift_demo.csv} 采用“长表”形式，每一行对应：

\begin{itemize}
    \item 一个对象 ID；
    \item 一个离散波长点；
    \item 一个归一化流量值；
    \item 数据角色（train/test/demo）；
    \item 一个粗略对象类别标签。
\end{itemize}

这个设计有几个优点：

\begin{itemize}
    \item 学生可以直接看到波长轴和流量轴的配对关系；
    \item 不需要二进制文件也能演示完整光谱思路；
    \item 既能说明吸收星、金属线主导星和发射源的差别，也能嵌入一个 redshift probe 示例。
\end{itemize}

\section{从长表重组到单条光谱}

本章的数据采用长表结构，这意味着同一个对象的多个波长点会分散在多行里。真正做谱线测量之前，先把它们按对象聚合成一条条小光谱，再按波长排序，通常是最稳妥的第一步。

\begin{examplebox}[按对象重组光谱]
\begin{lstlisting}[style=pythonstyle]
spectra = {}
for row in rows:
    sample_id = row["sample_id"]
    spectra.setdefault(sample_id, []).append(row)

for sample_rows in spectra.values():
    sample_rows.sort(key=lambda row: float(row["wavelength_rest_angstrom"]))
\end{lstlisting}
\end{examplebox}

这样做之后，你才能更自然地去比较连续谱、吸收线和发射线，并对候选谱线做等效宽度或红移测量。

\section{配套 notebook 做了什么}

本章 notebook 走的是一条典型而克制的路线：

\begin{itemize}
    \item 先按对象把长表重组成一条条小光谱；
    \item 比较几个对象的连续谱与谱线结构；
    \item 计算最小线深、发射峰和简化线比特征；
    \item 对一个 demo 对象做红移试算；
    \item 最后解释为什么归一化与采样一致性会影响后续分类或回归任务。
\end{itemize}

这条流程的重点不是一开始就做“真实大巡天级别”的光谱流水线，而是先把光谱分析的骨架讲清楚。

\section{常见错误}

\begin{warnbox}[光谱数据中的典型问题]
\begin{itemize}
    \item 只看流量数值，不看波长轴；
    \item 在不确认本征谱线身份的情况下直接估红移；
    \item 把没有归一化的整体亮度差异误读成谱线物理差异；
    \item 忽略采样不一致却直接比较不同光谱；
    \item 不考虑分辨率和信噪比，就过度解读微弱谱线；
    \item 只会“看见一个峰”，却说不清这个峰的本征波长和物理含义。
\end{itemize}
\end{warnbox}

\section{AI 助手如何帮助你}

在这一章，AI 助手适合帮助你：

\begin{itemize}
    \item 把光谱数据结构和代码整理得更清楚；
    \item 帮你根据谱线位置写出最小 redshift 计算步骤；
    \item 帮你总结连续谱、吸收线和发射线在图上的区别；
    \item 提醒你归一化、波长单位和重采样这些最容易被忽略的问题。
\end{itemize}

但某个谱线到底是否被正确识别、某个红移估计是否合理、归一化是否引入了新的偏差，仍然需要你自己做物理判断。

\section{练习}

\begin{exercisebox}[基础练习]
从本章教学数据里选一条光谱，找出它在 $4861\,\text{\AA}$ 和 $6563\,\text{\AA}$ 处是更像吸收线还是发射线，并说明你的判断依据。
\end{exercisebox}

\begin{exercisebox}[进阶练习]
已知一条谱线的本征波长为 $4861\,\text{\AA}$，若观测峰出现在 $4958\,\text{\AA}$，请计算红移并解释为什么这个数值意味着谱线向长波方向移动。
\end{exercisebox}

\begin{exercisebox}[开放练习]
为你自己的一个未来光谱项目写一份“最小光谱处理清单”，至少包含：波长单位、归一化方式、候选谱线、红移估计方法和采样一致性检查。
\end{exercisebox}

\section{本章小结}

光谱数据的核心，不是“很多波长点组成了一条曲线”，而是波长轴和流量结构一起记录了天体的物理状态。本章在最小教学光谱上补齐了从连续谱、分辨率、归一化、等效宽度到红移和速度判断的基本链条。只要这条链条真正建立起来，后面的分类、异常检测和更复杂的光谱机器学习任务就会更有物理抓手。
